\documentclass[11pt]{article}

\usepackage[a4paper,margin=1in]{geometry}
\usepackage{amsmath,amssymb,amsthm}
\usepackage{booktabs}
\usepackage{array}
\usepackage{xcolor}
\usepackage[colorlinks=true,linkcolor=blue!55!black,citecolor=blue!55!black,
            urlcolor=blue!55!black]{hyperref}

\allowdisplaybreaks
\numberwithin{equation}{section}

\newcommand{\E}{\mathbb{E}}
\newcommand{\R}{\mathbb{R}}
\newcommand{\one}{\mathbf{1}}
\newcommand{\code}[1]{\texttt{\small #1}}
\newcommand{\D}{\mathrm{D}}
\newcommand{\F}{\mathrm{F}}

\theoremstyle{definition}
\newtheorem{definition}{Definition}
\theoremstyle{plain}
\newtheorem{proposition}{Proposition}
\newtheorem{lemma}{Lemma}
\theoremstyle{remark}
\newtheorem{remark}{Remark}

\definecolor{notec}{RGB}{120,60,20}
\newcommand{\impl}[1]{\par\smallskip\noindent\textcolor{notec}{$\blacktriangleright$}\;\emph{#1}\par\smallskip}

\title{\vspace{-1.4cm}\bfseries Structural Equations of a Two-Country\\
Monetary Union with Sovereign Risk and\\ Constrained Financial Intermediaries\\[6pt]
\large A microfounded exposition of the model in \code{code/global/}}
\author{Model documentation}
\date{\today}

\begin{document}
\maketitle

\begin{abstract}
\noindent This document states the full structural environment of the two-country
heterogeneous-agent model of a monetary union solved in \code{code/global/}. The
economy consists of two countries, $\D$ (the risky sovereign; Greece) and $\F$ (the
core; Germany), each populated by incomplete-markets households, competitive
producers, capital producers with adjustment costs, monopolistically competitive
retailers, and Gertler--Karadi/Bocola financial intermediaries subject to an
occasionally-binding incentive-compatibility constraint. Only $\D$'s sovereign
debt carries default risk, and that risk is \emph{exogenous and priced}: following
Bocola (2016, \emph{JPE}), an increase in the priced one-period default probability
depresses long-duration bond prices, inflicts mark-to-market losses on intermediary
net worth, tightens the leverage constraint, widens lending spreads and contracts
output --- without default ever being realised along the baseline path. The
document derives each block from primitives, states every first-order condition,
gives the closed form of the leverage multiplier, defines the recursive competitive
equilibrium and the exact market-clearing system that the projection solver imposes,
and documents --- explicitly --- the approximations that separate the solved object
from the primitive environment.
\end{abstract}

\tableofcontents

%======================================================================
\section{Environment, timing and notation}
%======================================================================

\subsection{Overview}

Time is discrete and infinite, $t=0,1,2,\dots$, and a period is one quarter. The
world consists of two countries, indexed $X\in\{\D,\F\}$, forming a monetary union.
Each country produces a distinct traded good. Prices are fully flexible, so all
equilibrium objects below are \emph{real}; the ``monetary union'' is modelled at the
level at which it binds real allocations, namely a single union-wide funding market
for intermediary liabilities together with real interest parity
(Section~\ref{sec:money}).

Each country hosts four types of agent:
\begin{enumerate}
\item a continuum of \textbf{households} facing uninsurable idiosyncratic labour-income
      risk and a borrowing constraint, who save exclusively in bank deposits and supply
      labour with GHH preferences;
\item competitive \textbf{intermediate producers} who rent capital and hire labour, and
      who must pre-finance a fraction of the wage bill with intra-period bank credit
      (Neumeyer--Perri working capital);
\item \textbf{capital producers} subject to Jermann (1998) adjustment costs, and
      monopolistically competitive \textbf{retailers} with flexible prices;
\item a representative \textbf{financial intermediary} (``bank'') that holds the
      economy's productive capital, both sovereign bonds and the working-capital loan
      book, funds itself with household deposits, and is subject to an agency friction
      that generates an occasionally-binding leverage constraint.
\end{enumerate}
Each country also has a \textbf{government} issuing long-duration
Hatchondo--Martinez perpetuities and levying taxes according to a Bohn-type rule, and
the union has a \textbf{central bank} whose only instrument is a sovereign-bond
backstop of the Transmission Protection Instrument (TPI/OMT) type.

\subsection{Timing convention}

Timing conventions are load-bearing in this model and are stated once, here.

\begin{enumerate}
\item \textbf{Capital is predetermined} (Bocola, 2016, eq.~6). The capital stock that
      produces at $t$ was purchased and priced by banks at $t-1$. Denote by $K_{X,t}$
      the stock \emph{producing} at $t$; the bank chooses $K_{X,t+1}$ at $t$. Impact
      responses of output therefore run through hours alone.
\item \textbf{The deposit rate is predetermined}. The gross rate paid on deposits held
      from $t$ to $t+1$, $1+r_{X,t}$, is set at $t$. It is the return that appears in
      the household Euler equation formed at $t$ and in the bank's funding cost at
      $t+1$.
\item \textbf{Default is realised at the beginning of a period}, before payoffs are
      collected: a haircut applies to coupons and to the continuation value of the
      whole outstanding claim.
\item \textbf{Working-capital loans are intra-period}: extended at $t$, repaid at $t$
      with interest, and settled inside the bank's carried obligation, so they never
      appear as a separate state.
\end{enumerate}

\subsection{Notation}

Table~\ref{tab:notation} fixes symbols. Two conventions deserve emphasis. First,
$p_t$ is the terms of trade measured as \emph{units of the $\D$-good per unit of the
$\F$-good}, so $p_t\uparrow$ is a real depreciation for $\D$. Second, $P_{X,t}$
denotes the bank's \emph{gross deposit obligation} carried into $t$ (a state
variable, following Bocola's notation), while $P^{c}_{X,t}$ denotes the CES
\emph{consumption price index}. They are distinct objects.

\begin{table}[htbp]
\centering
\small
\caption{Notation. Country index $X\in\{\D,\F\}$ suppressed where unambiguous.}
\label{tab:notation}
\begin{tabular}{@{}ll@{\qquad}ll@{}}
\toprule
Symbol & Meaning & Symbol & Meaning \\
\midrule
$C_{X,t}$        & aggregate consumption (basket units) & $n_{X,t}$      & bank net worth (end of $t$)\\
$N_{X,t}$        & aggregate hours                      & $n^{g}_{X,t}$  & gross bank wealth (start of $t$)\\
$a_{it}$         & individual deposit holdings          & $\mathcal{A}_{X,t}$ & bank total assets (market value)\\
$A_{X,t}$        & aggregate real deposits              & $P_{X,t}$      & gross deposit obligation (state)\\
$e_{it}$         & idiosyncratic labour efficiency      & $\varphi_{X,t}$ & marginal value of bank net worth\\
$Y_{X,t}$        & gross output                         & $\theta_{X,t}$  & bank leverage $\mathcal{A}/n$\\
$K_{X,t}$        & capital \emph{producing} at $t$       & $\mu_{X,t}$     & IC multiplier (rescaled)\\
$I_{X,t}$        & investment                           & $\lambda$       & divertable asset share\\
$Q_{X,t}$        & Tobin's $q$ (capital price)          & $f$             & banker exit/payout share\\
$w_{X,t}$        & real wage (post-wedge)               & $\omega$        & entrant endowment share\\
$r_{X,t}$        & deposit rate set at $t$              & $\Lambda_{t,t+1}$ & household SDF\\
$r^{k}_{X,t}$    & realised return on capital           & $\Omega_{t,t+1}$  & bank's augmented SDF\\
$r^{wc}_{X,t}$   & working-capital lending rate         & $\zeta$         & working-capital share of wage bill\\
$L_{X,t}$        & working-capital loans                & $\delta_b$      & bond amortisation rate\\
$p_{t}$          & terms of trade ($\D$ per $\F$)       & $B_{t}$         & $\D$-sovereign stock (face)\\
$P^{c}_{X,t}$    & CES consumption price index          & $q^{X}_{t}$     & sovereign bond price\\
$IM_{X,t}$       & imports                              & $s_{t}$         & sovereign-risk factor\\
$NX_{X,t}$       & net exports (own-good units)         & $\pi^{d}_{t}$   & priced default probability\\
$Z_{X,t}$        & total factor productivity            & $d_{t}$         & realised default indicator\\
$T^{\tau}_{X,t}$ & lump-sum taxes                       & $h_{t}$         & survival factor $1-d_t(1-\varrho)$\\
$\mathrm{Div}_{X,t}$ & dividends to households          & $\varrho$       & recovery rate\\
$\mathcal{S}_t$  & aggregate state vector               & $\varkappa_{\rm tpi}$ & priced TPI activation prob.\\
\bottomrule
\end{tabular}
\end{table}

%======================================================================
\section{Households}
%======================================================================

\subsection{Preferences and the individual problem}

Country $X$ is populated by a unit continuum of infinitely-lived households indexed
by $i$. Preferences are of the Greenwood--Hercowitz--Huffman (GHH) form: the
consumption--labour composite
\begin{equation}
x_{it}\;\equiv\;c_{it}-v(n_{it}),
\qquad
v(n)\;=\;\chi_X\,\frac{n^{\,1+1/\nu_X}}{1+1/\nu_X},
\label{eq:ghh-composite}
\end{equation}
enters a CRRA felicity function, so that household $i$ solves
\begin{equation}
\max_{\{c_{it},\,n_{it},\,a_{i,t+1}\}}\;
\E_0\sum_{t=0}^{\infty}\beta_X^{\,t}\,
\frac{\bigl(c_{it}-v(n_{it})\bigr)^{1-\sigma_X}-1}{1-\sigma_X},
\label{eq:hh-objective}
\end{equation}
subject to the flow budget constraint, expressed in units of the country's
consumption basket,
\begin{equation}
c_{it}+a_{i,t+1}
\;=\;
\bigl(1+r^{c}_{X,t}\bigr)\,a_{it}
\;+\;\frac{w_{X,t}}{P^{c}_{X,t}}\,e_{it}\,n_{it}
\;+\;\frac{\mathrm{Div}_{X,t}-T^{\tau}_{X,t}}{P^{c}_{X,t}},
\label{eq:hh-budget}
\end{equation}
and to the no-borrowing constraint
\begin{equation}
a_{i,t+1}\;\ge\;\underline{a}_X \;=\;0 .
\label{eq:hh-borrowing}
\end{equation}
Here $a_{it}$ are real deposits carried into $t$, the only asset available to
households, and
\begin{equation}
1+r^{c}_{X,t}\;\equiv\;\bigl(1+r_{X,t-1}\bigr)\,\frac{P^{c}_{X,t-1}}{P^{c}_{X,t}}
\label{eq:hh-realreturn}
\end{equation}
is the realised real return on the deposit contract signed at $t-1$ --- the
predetermined-rate convention of Section~1.2. Dividends $\mathrm{Div}_{X,t}$ are
the consolidated payout of retailers, capital producers and exiting bankers
(equation~\eqref{eq:dividends}); $T^{\tau}_{X,t}$ are lump-sum taxes.

Idiosyncratic labour efficiency $e_{it}$ follows a stationary AR(1) in logs,
\begin{equation}
\log e_{i,t+1}=\rho_e\log e_{it}+\epsilon_{i,t+1},
\qquad \epsilon\sim N(0,\sigma_e^2),
\label{eq:income-ar1}
\end{equation}
discretised by the Rouwenhorst (1995) method into an $n_e$-state Markov chain
$(\mathsf{e},\Pi)$ with grid normalised so that $\sum_e \pi^{\ast}(e)\,\mathsf{e}(e)=1$,
where $\pi^{\ast}$ is the stationary distribution of $\Pi$. Idiosyncratic risk is
uninsurable: markets are incomplete and the only self-insurance vehicle is the
deposit account, subject to~\eqref{eq:hh-borrowing}.

\subsection{First-order conditions}

Let $\eta_{it}\ge0$ be the multiplier on the borrowing constraint. The
intertemporal and intratemporal optimality conditions are
\begin{align}
x_{it}^{-\sigma_X}
&\;=\;\beta_X\,\bigl(1+r_{X,t}\bigr)\,
   \E_t\!\left[\Bigl(\tfrac{P^{c}_{X,t}}{P^{c}_{X,t+1}}\Bigr)\,x_{i,t+1}^{-\sigma_X}\right]
   \;+\;\eta_{it},
\qquad
\eta_{it}\bigl(a_{i,t+1}-\underline{a}_X\bigr)=0,
\label{eq:hh-euler}\\[4pt]
\chi_X\,n_{it}^{1/\nu_X}
&\;=\;\frac{w_{X,t}}{P^{c}_{X,t}}\;e_{it}.
\label{eq:hh-labour}
\end{align}
Equation~\eqref{eq:hh-labour} is the defining analytical convenience of GHH
preferences: the marginal rate of substitution between consumption and leisure is
independent of $c_{it}$, hence of wealth. Labour supply per efficiency unit is
therefore \emph{identical across households}, and aggregation over the wealth
distribution is exact. Writing $N_{X,t}\equiv\int e_{it}n_{it}\,d\Phi_t$ for
efficiency-weighted aggregate hours and using $\int e\,d\Phi=1$, the aggregate
labour-supply schedule is
\begin{equation}
\boxed{\;\chi_X\,N_{X,t}^{1/\nu_X}\;=\;\frac{w_{X,t}}{P^{c}_{X,t}}\;}
\label{eq:agg-labour}
\end{equation}
which is the object the equilibrium system imposes (residual 3--4 in
Section~\ref{sec:system}). The same property implies that the intertemporal margin
depends on the composite $x_{it}$ only, so \eqref{eq:hh-euler} can be solved by the
endogenous grid method on $x$.

\impl{GHH removes the wealth effect on labour supply. This is what allows a
distribution-free aggregate labour block to coexist with a genuinely
heterogeneous-agent savings block, but it also means that redistribution has no
direct effect on hours --- a property that matters for how the working-capital
financing income is routed (Section~\ref{sec:wc}).}

\subsection{Distribution and aggregation}

Let $\Phi_t(a,e)$ denote the joint distribution over deposits and efficiency, and
$g^{a}_{X,t}(a,e)$, $g^{c}_{X,t}(a,e)$ the optimal savings and consumption policies
implied by \eqref{eq:hh-euler}--\eqref{eq:hh-labour}. The distribution evolves under
the standard forward operator
\begin{equation}
\Phi_{t+1}(a',e')\;=\;\sum_{e}\Pi(e,e')\int
   \one\bigl\{g^{a}_{X,t}(a,e)=a'\bigr\}\,d\Phi_t(a,e),
\label{eq:kfe}
\end{equation}
implemented numerically by the Young (2010) lottery: an off-grid choice
$g^{a}\in[a_j,a_{j+1}]$ is split between the adjacent nodes with weights
$1-\omega_j$ and $\omega_j=(g^{a}-a_j)/(a_{j+1}-a_j)$, which preserves the first
moment exactly. Aggregates are
\begin{equation}
A_{X,t+1}=\int g^{a}_{X,t}\,d\Phi_t,
\qquad
C_{X,t}=\int g^{c}_{X,t}\,d\Phi_t .
\label{eq:hh-aggregates}
\end{equation}

\subsection{The aggregate closure used in the recursive solution}
\label{sec:hh-closure}

The heterogeneous-agent block \eqref{eq:hh-objective}--\eqref{eq:kfe} is solved
exactly at the steady state (it pins the discount factors $\beta_X$ from
deposit-market clearing) and is used for the distributional/welfare overlays. In the
recursive global solution, however, the household side is closed at the
\emph{aggregate} level by a representative-agent Euler equation --- the first rung of
a Krusell--Smith ladder, in which the cross-sectional distribution is carried by its
mean alone. Three equations define this closure.

\paragraph{(i) Quantity clearing of the deposit market.} Aggregate real household
deposits equal aggregate bank funding, by construction:
\begin{equation}
A_{X,t}\;=\;\frac{\mathrm{dep}_{X,t}}{P^{c}_{X,t}} .
\label{eq:dep-quantity}
\end{equation}

\paragraph{(ii) The aggregate budget constraint.} Consumption is the residual of the
household budget, where the gross deposit \emph{claim} carried into $t$ is the bank's
gross deposit \emph{obligation} state $P_{X,t}$:
\begin{equation}
C_{X,t}
\;=\;\underbrace{\frac{P_{X,t}}{P^{c}_{X,t}}}_{\text{carried wealth}}
\;+\;\underbrace{\frac{w_{X,t}N_{X,t}+\mathrm{Div}_{X,t}-T^{\tau}_{X,t}}{P^{c}_{X,t}}
     +\bar{T}_X}_{\;\equiv\;\text{inc}_{X,t}}
\;-\;A_{X,t}.
\label{eq:hh-budget-agg}
\end{equation}

\paragraph{(iii) The aggregate deposit Euler equation.} With $x_{X,t}=C_{X,t}-v(N_{X,t})$,
\begin{equation}
\boxed{\;
x_{X,t}^{-\sigma_X}
=\tilde\beta_X\,
\E_t\!\left[\bigl(1+r^{c}_{X,t+1}\bigr)\,x_{X,t+1}^{-\sigma_X}\right],
\qquad
1+r^{c}_{X,t+1}=\bigl(1+r_{X,t}\bigr)\frac{P^{c}_{X,t}}{P^{c}_{X,t+1}} \;}
\label{eq:agg-euler}
\end{equation}
with $\tilde\beta_X=1/(1+\bar r_X)$, where $\bar r_X$ is the steady-state deposit
rate. Equation~\eqref{eq:agg-euler} is the condition that determines the deposit
\emph{rate} $r_{\D,t}$ in equilibrium (residual 5 of Section~\ref{sec:system}):
quantities clear by \eqref{eq:dep-quantity}, and the rate is whatever makes
households willing to hold exactly what banks fund.

\begin{remark}[Two disclosed departures from the primitive environment]
\label{rem:closure}
The representative-agent stand-in is not the aggregate of
\eqref{eq:hh-objective}: (a) $\tilde\beta_X=1/(1+\bar r_X)$ exceeds the
heterogeneous-agent $\beta_X$ recovered at the steady state, because precautionary
saving in the incomplete-markets block requires $\beta_X(1+\bar r_X)<1$; and (b) the
constant $\bar T_X$ in \eqref{eq:hh-budget-agg} is an accounting anchor, calibrated
so that the steady state is an exact rest point of the recursive system. It absorbs
the working-capital repayment leg, which is netted out of the state $P_{X,t}$ (see
equation~\eqref{eq:P-state}) but accrues to households. Both are stated here rather
than buried: the recursive solution is exact in the aggregate but treats the wealth
\emph{distribution} as a fixed shape.
\end{remark}

%======================================================================
\section{Production, trade and capital}
%======================================================================

\subsection{Final goods and the terms of trade}

The consumption basket of country $\D$ aggregates home and foreign goods with an
Armington/CES aggregator,
\begin{equation}
C_{\D,t}=\Bigl[\varpi^{1/\eta}c_{\D\D,t}^{\frac{\eta-1}{\eta}}
+(1-\varpi)^{1/\eta}c_{\F\D,t}^{\frac{\eta-1}{\eta}}\Bigr]^{\frac{\eta}{\eta-1}},
\label{eq:ces-aggregator}
\end{equation}
with home bias $\varpi$ and trade elasticity $\eta$. Normalising the price of the
$\D$-good to one and letting $p_t$ be the relative price of the $\F$-good, cost
minimisation gives the price index and the import demand schedule
\begin{align}
P^{c}_{\D,t}&=\Bigl[\varpi+(1-\varpi)\,p_t^{\,1-\eta}\Bigr]^{\frac{1}{1-\eta}},
&
IM_{\D,t}&=(1-\varpi)\Bigl(\frac{P^{c}_{\D,t}}{p_t}\Bigr)^{\eta}C_{\D,t},
\label{eq:ces-D}\\[3pt]
P^{c}_{\F,t}&=\Bigl[\varpi+(1-\varpi)\,p_t^{\,\eta-1}\Bigr]^{\frac{1}{1-\eta}},
&
IM_{\F,t}&=(1-\varpi)\bigl(P^{c}_{\F,t}\,p_t\bigr)^{\eta}C_{\F,t},
\label{eq:ces-F}
\end{align}
where $P^{c}_{\F,t}$ is expressed in $\F$-good units and the $\D$-good costs $1/p_t$
there. Net exports in own-good units are
\begin{equation}
NX_{\D,t}=IM_{\F,t}-p_t\,IM_{\D,t},
\qquad
NX_{\F,t}=IM_{\D,t}-\frac{IM_{\F,t}}{p_t}.
\label{eq:nx}
\end{equation}

\subsection{Retailers and marginal cost}

A unit continuum of retailers buys the homogeneous intermediate good, differentiates
it costlessly, and sells under monopolistic competition to a CES final-goods
aggregator with elasticity $\epsilon_X$. Prices are fully flexible, so every retailer
charges the constant markup $\epsilon_X/(\epsilon_X-1)$ over marginal cost and real
marginal cost is constant:
\begin{equation}
mc_X=\frac{\epsilon_X-1}{\epsilon_X}.
\label{eq:mc}
\end{equation}
Retail profits $(1-mc_X)Y_{X,t}$ are rebated lump-sum to households.

\subsection{Intermediate producers and the working-capital wedge}
\label{sec:wc}

The intermediate producer operates a Cobb--Douglas technology in the
\emph{predetermined} capital stock and hours,
\begin{equation}
Y_{X,t}=Z_{X,t}\,K_{X,t}^{\alpha_X}\,N_{X,t}^{1-\alpha_X},
\label{eq:production}
\end{equation}
where $K_{X,t}$ was purchased by banks at $t-1$. Following Bocola's own open-economy
extension (his \S V.C), the firm must pre-finance a fraction $\zeta_X$ of its wage
bill with intra-period bank credit at the gross rate $1+r^{wc}_{X,t}$. The static
problem is
\begin{equation}
\max_{N_{X,t}}\;\;
mc_X\,Z_{X,t}K_{X,t}^{\alpha_X}N_{X,t}^{1-\alpha_X}
\;-\;\bigl(1+\zeta_X r^{wc}_{X,t}\bigr)\,w_{X,t}N_{X,t}
\;-\;mpk_{X,t}K_{X,t},
\label{eq:firm-problem}
\end{equation}
delivering the factor-demand conditions
\begin{align}
w_{X,t}&=\frac{mc_X(1-\alpha_X)\,Y_{X,t}/N_{X,t}}{1+\zeta_X\,r^{wc}_{X,t}},
\label{eq:labour-demand}\\[3pt]
mpk_{X,t}&=mc_X\,\alpha_X\,\frac{Y_{X,t}}{K_{X,t}},
\label{eq:mpk}
\end{align}
and a loan demand of
\begin{equation}
L_{X,t}=\zeta_X\,w_{X,t}N_{X,t}.
\label{eq:loan-demand}
\end{equation}
Note that $mpk_{X,t}$ in \eqref{eq:mpk} is the marginal product of the
\emph{bank-held vintage}: firms rent exactly the stock banks financed one period
earlier.

\impl{Equation~\eqref{eq:labour-demand} is the \textbf{only} channel through which
financial spreads reach output on impact. Setting $\zeta_X=0$ nests the model
without it exactly, and doing so collapses the output response to a sovereign-risk
shock to approximately zero even when bond prices fall by 30\%. The interest cost
$\zeta_X r^{wc}wN$ is treated as a real resource cost in the baseline; a switch
($\varsigma^{\rm reb}\in\{0,1\}$ in equation~\eqref{eq:dividends}) can instead rebate
it to households, in which case, under GHH, the spread becomes a pure intra-period
transfer with no offsetting wealth effect and hours \emph{rise} with the spread.}

\subsection{\texorpdfstring{Capital producers, Tobin's $q$, and the return on capital}
                          {Capital producers, Tobin's q, and the return on capital}}

Capital producers convert the final good into installed capital subject to Jermann
(1998) adjustment costs. With $\iota_{X,t}\equiv I_{X,t}/K_{X,t}$, the law of motion is
\begin{equation}
K_{X,t+1}=(1-\delta_X)K_{X,t}+\Phi_X(\iota_{X,t})\,K_{X,t},
\qquad
\Phi_X(\iota)=\gamma_{0,X}\,\iota^{\,1-\xi_X}+\gamma_{1,X},
\label{eq:capital-lom}
\end{equation}
with the coefficients pinned so that the adjustment cost is inactive at the steady
state ($\Phi_X(\delta_X)=\delta_X$, $\Phi_X'(\delta_X)=1$):
\begin{equation}
\gamma_{0,X}=\frac{\delta_X^{\;\xi_X}}{1-\xi_X},
\qquad
\gamma_{1,X}=-\frac{\delta_X\,\xi_X}{1-\xi_X}.
\label{eq:jermann-coeffs}
\end{equation}
The producer solves $\max_{\iota}\;Q_{X,t}\Phi_X(\iota)K_{X,t}-\iota K_{X,t}$, whose
first-order condition delivers marginal Tobin's $q$ and, inverting
\eqref{eq:capital-lom}, the investment rate:
\begin{equation}
Q_{X,t}=\frac{1}{\Phi_X'(\iota_{X,t})}
      =\frac{\iota_{X,t}^{\;\xi_X}}{\gamma_{0,X}(1-\xi_X)},
\qquad
\iota_{X,t}=\left[\frac{K_{X,t+1}/K_{X,t}-(1-\delta_X)-\gamma_{1,X}}
                       {\gamma_{0,X}}\right]^{\frac{1}{1-\xi_X}} .
\label{eq:tobins-q}
\end{equation}
Capital producers' rents, $\Pi^{k}_{X,t}=Q_{X,t}\bigl[K_{X,t+1}-(1-\delta_X)K_{X,t}\bigr]-I_{X,t}$,
are rebated to households. The parameter $\xi_X$ is the elasticity of $Q$ with respect
to $I/K$.

The bank holds the capital, so the \emph{realised} gross return between $t$ and $t+1$
on a unit of capital purchased at $t$ is
\begin{equation}
\boxed{\;1+r^{k}_{X,t+1}=\frac{mpk_{X,t+1}+(1-\delta_X)\,Q_{X,t+1}}{Q_{X,t}}\;}
\label{eq:rk}
\end{equation}
which is the object that appears in the intermediary's capital Euler equation.

Total factor productivity in $\D$ is carried as a state and follows a deterministic
(perfect-foresight) AR(1),
\begin{equation}
\log Z_{\D,t+1}=(1-\rho_z)\log \bar Z_{\D}+\rho_z\log Z_{\D,t},
\label{eq:tfp}
\end{equation}
while $Z_{\F,t}=\bar Z_{\F}$ throughout: country asymmetries enter through shocks
only.

%======================================================================
\section{Financial intermediaries}
\label{sec:banks}
%======================================================================

This is the core of the model. Each country hosts a representative bank that
intermediates \emph{all} of the economy's productive capital, the sovereign bonds
held domestically, a cross-border sovereign position, and the working-capital loan
book. Its balance sheet is the transmission mechanism from sovereign risk to real
activity.

\subsection{The balance sheet}

At the end of period $t$ the bank in country $\D$ holds capital $K_{\D,t+1}$ at
price $Q_{\D,t}$, home sovereign bonds $b^{\D}_{\D,t+1}$ at price $q^{\D}_t$,
foreign sovereign bonds $b^{\F}_{\D,t+1}$ at price $q^{\F}_t$ converted into
$\D$-goods at the terms of trade $p_t$, and working-capital loans $L_{\D,t}$. These
are funded by net worth $n_{\D,t}$ and deposits $\mathrm{dep}_{\D,t}$:
\begin{equation}
\underbrace{Q_{\D,t}K_{\D,t+1}
   +q^{\D}_{t}b^{\D}_{\D,t+1}
   +p_t\,q^{\F}_{t}b^{\F}_{\D,t+1}
   +L_{\D,t}}_{\textstyle \equiv\;\mathcal{A}_{\D,t}\ \ (\text{total assets})}
\;=\;n_{\D,t}+\mathrm{dep}_{\D,t}.
\label{eq:balance-sheet}
\end{equation}
The mirror-image expression for $\F$ converts its cross-border leg by $1/p_t$. Bank
leverage is $\theta_{X,t}\equiv\mathcal{A}_{X,t}/n_{X,t}$.

\subsection{Gross wealth and the deposit-obligation state}

Let $h_t\equiv1-d_t(1-\varrho_\D)$ denote the survival factor on $\D$-sovereign
claims (Section~\ref{sec:sovrisk}). At the start of $t+1$ the bank collects
\begin{equation}
\mathcal{X}_{\D,t+1}
=\bigl[mpk_{\D,t+1}+(1-\delta_\D)Q_{\D,t+1}\bigr]K_{\D,t+1}
+\Xi^{\D}_{t+1}\,b^{\D}_{\D,t+1}
+p_{t+1}\,\Xi^{\F}_{t+1}\,b^{\F}_{\D,t+1},
\label{eq:asset-payoff}
\end{equation}
where the per-unit bond payoffs are (Section~\ref{sec:sovrisk})
\begin{equation}
\Xi^{\D}_{t+1}=h_{t+1}\bigl[\delta_b+(1-\delta_b)q^{\D}_{t+1}\bigr],
\qquad
\Xi^{\F}_{t+1}=\delta_b+(1-\delta_b)q^{\F}_{t+1}.
\label{eq:bond-payoff}
\end{equation}
Against this it owes its predetermined obligation. Because the working-capital loan
is extended and repaid \emph{within} the period at the rate locked at $t$, the
obligation the bank carries forward is net of the loan receivable:
\begin{equation}
\boxed{\;P_{X,t+1}\;=\;\bigl(1+r_{X,t}\bigr)\,\mathrm{dep}_{X,t}
\;-\;\bigl(1+r^{wc}_{X,t}\bigr)L_{X,t}\;}
\label{eq:P-state}
\end{equation}
so that gross bank wealth at the start of $t+1$ is simply
\begin{equation}
n^{g}_{X,t+1}=\mathcal{X}_{X,t+1}-P_{X,t+1}.
\label{eq:ng}
\end{equation}
This is exactly the algebraic device that makes $P_{X}$ a \emph{sufficient} state for
the bank's liability side: neither deposits, nor the loan book, nor the loan rate
need to be carried separately.

Substituting \eqref{eq:balance-sheet} into \eqref{eq:ng} and using the return
definitions gives the excess-return representation used throughout:
\begin{equation}
n^{g}_{X,t+1}
=\sum_{j\in\mathcal{J}}\bigl(R_{j,t+1}-R_{X,t}\bigr)\,a_{j,t}
\;+\;R_{X,t}\,n_{X,t},
\qquad R_{X,t}\equiv1+r_{X,t},
\label{eq:ng-excess}
\end{equation}
where $\mathcal{J}=\{K,\,b^{\D},\,b^{\F},\,L\}$ indexes asset classes, $a_{j,t}$ is
the market value of position $j$, and $R_{j,t+1}$ its gross return. The
working-capital leg is riskless as of $t$: $R_{L,t+1}=1+r^{wc}_{X,t}$.

\subsection{The banker's problem and the franchise value}

Each period a fraction $f$ of bankers exits and pays its net worth out as dividends;
the remaining $1-f$ continue. Entrants receive a startup transfer equal to a fraction
$\omega$ of the aggregate asset base. A continuing banker maximises the expected
discounted value of terminal net worth, discounting with the household stochastic
discount factor $\Lambda_{t,t+1}$ of the country it operates in. Writing the value
function of a banker with net worth $n$ as $V_t(n)$, the Bellman equation is
\begin{equation}
V_t(n_{X,t})
=\max_{\{a_{j,t}\}}\;
\E_t\Bigl\{\Lambda_{X,t,t+1}\bigl[
   f\,n^{g}_{X,t+1}+(1-f)\,V_{t+1}\bigl(n^{g}_{X,t+1}\bigr)\bigr]\Bigr\}.
\label{eq:bank-bellman}
\end{equation}
Because \eqref{eq:ng-excess} is linear in $n$ and the constraint below is linear in
assets, the value function is linear, $V_t(n)=\varphi_{X,t}\,n$, where
$\varphi_{X,t}$ is the \textbf{marginal value of one unit of intermediary net
worth} (the franchise value). Defining the bank's \textbf{augmented stochastic
discount factor}
\begin{equation}
\boxed{\;
\Omega_{X,t,t+1}\;\equiv\;\Lambda_{X,t,t+1}\Bigl[\,f+(1-f)\,\varphi_{X,t+1}\Bigr]\;}
\label{eq:omega}
\end{equation}
the objective collapses to $\E_t\bigl[\Omega_{X,t,t+1}\,n^{g}_{X,t+1}\bigr]$. The
weight on the franchise value is the \emph{survival} probability $1-f$; the weight
$f$ is the exit/payout share.

\subsection{The agency friction and the incentive-compatibility constraint}

After raising deposits, the banker can abscond with a fraction $\lambda$ of assets,
in which case depositors recover the rest and the bank is liquidated. Deposits are
forthcoming only if the franchise is worth at least as much as the divertable
proceeds:
\begin{equation}
\boxed{\;\varphi_{X,t}\,n_{X,t}\;\ge\;\lambda_X\,\mathcal{A}_{X,t}\;}
\label{eq:IC}
\end{equation}
Following Bocola (2016, eq.~3), \emph{all} asset classes carry the same divertability
$\lambda_X$ --- capital, home bonds, cross-border bonds and working-capital loans
alike. Equation~\eqref{eq:IC} is an endogenous leverage ceiling:
$\theta_{X,t}\le\varphi_{X,t}/\lambda_X$.

\begin{remark}[Why a single $\lambda$]
Allowing $\lambda_K\ne\lambda_{b^{\D}}$ re-opens a portfolio-substitution margin
along which a fall in sovereign bond prices \emph{relaxes} the constraint (the bank
shifts into the now-cheaper divertability class), which makes sovereign risk
expansionary. The single-$\lambda$ restriction is therefore not a simplification but
a structural choice inherited from Bocola.
\end{remark}

\subsection{First-order conditions and the closed-form multiplier}

Attach the multiplier $\tilde\mu_{X,t}\ge0$ to \eqref{eq:IC} and write the
Lagrangian of the banker's problem as
\begin{equation}
\mathcal{L}
=\bigl(1+\tilde\mu_{X,t}\bigr)\,
  \E_t\bigl[\Omega_{X,t,t+1}\,n^{g}_{X,t+1}\bigr]
  \;-\;\tilde\mu_{X,t}\,\lambda_X\,\mathcal{A}_{X,t},
\label{eq:bank-lagrangian}
\end{equation}
using $V_t(n)=\E_t[\Omega\,n^{g}_{t+1}]$ on the left of \eqref{eq:IC}. It is
convenient to rescale the multiplier onto the unit interval,
\begin{equation}
\mu_{X,t}\;\equiv\;\frac{\tilde\mu_{X,t}}{1+\tilde\mu_{X,t}}\;\in\;[0,1).
\label{eq:mu-rescale}
\end{equation}
Differentiating \eqref{eq:bank-lagrangian} with respect to each position $a_{j,t}$
and using \eqref{eq:ng-excess} gives the \textbf{unified asset-pricing condition}
\begin{equation}
\boxed{\;\E_t\Bigl[\Omega_{X,t,t+1}\bigl(R_{j,t+1}-R_{X,t}\bigr)\Bigr]
\;=\;\lambda_X\,\mu_{X,t},
\qquad j\in\{K,\;b^{\D},\;b^{\F},\;L\}\;}
\label{eq:foc-general}
\end{equation}
together with the Karush--Kuhn--Tucker complementarity
\begin{equation}
\mu_{X,t}\ge0,
\qquad
\mathrm{slack}_{X,t}\equiv\varphi_{X,t}n_{X,t}-\lambda_X\mathcal{A}_{X,t}\ge0,
\qquad
\mu_{X,t}\cdot\mathrm{slack}_{X,t}=0 .
\label{eq:kkt}
\end{equation}
Every asset class earns the \emph{same} shadow excess return $\lambda_X\mu_{X,t}$:
the constraint prices scarce intermediary equity, not asset-specific risk.

\begin{proposition}[Franchise value and the closed-form multiplier; Bocola 2016, Prop.~1]
\label{prop:mu}
At an optimum the franchise value satisfies the recursion
\begin{equation}
\varphi_{X,t}\;=\;\frac{\E_t\bigl[\Omega_{X,t,t+1}\bigr]\,R_{X,t}}{1-\mu_{X,t}},
\label{eq:phi-recursion}
\end{equation}
and, whenever the incentive constraint binds, the multiplier admits the closed form
\begin{equation}
\mu_{X,t}
=\max\left\{\,1-\frac{\E_t\bigl[\Omega_{X,t,t+1}\bigr]\,R_{X,t}\;n_{X,t}}
                     {\lambda_X\,\mathcal{A}_{X,t}}\;,\;0\right\}.
\label{eq:mu-closed}
\end{equation}
\end{proposition}

\begin{proof}
Evaluate the objective at the optimum using \eqref{eq:ng-excess} and
\eqref{eq:foc-general}:
\begin{align*}
\varphi_{X,t}n_{X,t}
&=\E_t\bigl[\Omega\,n^{g}_{t+1}\bigr]
 =\sum_j \E_t\bigl[\Omega(R_{j,t+1}-R_{X,t})\bigr]a_{j,t}
  +\E_t[\Omega]\,R_{X,t}n_{X,t}\\
&=\mu_{X,t}\,\lambda_X\mathcal{A}_{X,t}+\E_t[\Omega]\,R_{X,t}n_{X,t}.
\end{align*}
If the constraint binds, $\lambda_X\mathcal{A}_{X,t}=\varphi_{X,t}n_{X,t}$, so
$\varphi_{X,t}n_{X,t}(1-\mu_{X,t})=\E_t[\Omega]R_{X,t}n_{X,t}$, which is
\eqref{eq:phi-recursion}; substituting \eqref{eq:phi-recursion} back into the binding
constraint and solving for $\mu$ gives \eqref{eq:mu-closed}. If the constraint is
slack, $\mu_{X,t}=0$ and \eqref{eq:phi-recursion} still holds, because the excess-return
terms vanish by \eqref{eq:foc-general}. The $\max\{\cdot,0\}$ operator is exactly the
complementarity \eqref{eq:kkt}.
\end{proof}

\impl{Equation~\eqref{eq:mu-closed} is the analytical heart of the transmission
mechanism. A fall in sovereign bond prices reduces $n_{X,t}$ through
\eqref{eq:ng}, which \emph{mechanically} raises $\mu_{X,t}$: scarcer equity, a
tighter constraint, higher required excess returns on every asset class, and --- via
\eqref{eq:rwc} --- a higher cost of working capital and lower hours. The premium is
endogenous and requires no ad hoc spread rule.}

\subsection{The three pricing conditions in explicit form}

Specialising \eqref{eq:foc-general} to each class yields the equations imposed
in equilibrium.

\paragraph{Capital.} With $R_{K,t+1}=1+r^{k}_{X,t+1}$ from \eqref{eq:rk}:
\begin{equation}
\E_t\Bigl[\Omega_{X,t,t+1}\bigl(r^{k}_{X,t+1}-r_{X,t}\bigr)\Bigr]
=\lambda_X\,\mu_{X,t}.
\label{eq:capital-euler}
\end{equation}
This is the equation that pins down the capital stock $K_{X,t+1}$ carried into the
next period (residuals 1--2 of Section~\ref{sec:system}).

\paragraph{Sovereign bonds.} With $R_{b,t+1}=\Xi_{t+1}/q_t$, condition
\eqref{eq:foc-general} rearranges into an explicit pricing formula,
\begin{align}
q^{\D}_{t}&=\frac{\E_t\bigl[\Omega_{\D,t,t+1}\,\Xi^{\D}_{t+1}\bigr]}
                 {\E_t\bigl[\Omega_{\D,t,t+1}\bigr]R_{\D,t}+\lambda_\D\,\mu_{\D,t}},
&
q^{\F}_{t}&=\frac{\E_t\bigl[\Omega_{\F,t,t+1}\,\Xi^{\F}_{t+1}\bigr]}
                 {\E_t\bigl[\Omega_{\F,t,t+1}\bigr]R_{\F,t}+\lambda_\F\,\mu_{\F,t}} .
\label{eq:bond-pricing}
\end{align}
The denominator makes the \emph{liquidity} (constraint) component of the yield
explicit: even a default-free bond trades below its risk-neutral present value by the
factor $\lambda\mu$. The numerator carries the default risk.

\paragraph{Working capital.} The loan rate is set at $t$ and is riskless, so
$\E_t[\Omega](1+r^{wc}_{X,t}-R_{X,t})=\lambda_X\mu_{X,t}$, i.e.
\begin{equation}
\boxed{\;
r^{wc}_{X,t}\;=\;r_{X,t}\;+\;\underbrace{\frac{\lambda_X\,\mu_{X,t}}
                                               {\E_t\bigl[\Omega_{X,t,t+1}\bigr]}}
                                        _{\text{lending spread}}\;}
\label{eq:rwc}
\end{equation}
Equation~\eqref{eq:rwc} is the model's \textbf{credit spread}. It is the single
quantity that links the financial block to the real block, through
\eqref{eq:labour-demand}.

\subsection{Net worth accumulation, entry and exit}

Aggregating over bankers, net worth carried out of period $t$ is the retained wealth
of surviving bankers plus the endowment of entrants:
\begin{equation}
\boxed{\;n_{X,t}\;=\;(1-f)\,n^{g}_{X,t}\;+\;\omega_X\,\mathcal{A}_{X,t}\;}
\label{eq:nw-lom}
\end{equation}
with $n^{g}_{X,t}=\mathcal{X}_{X,t}-P_{X,t}$ from \eqref{eq:ng}, and the corresponding
net payout to households is
\begin{equation}
\mathrm{div}^{\rm bank}_{X,t}=f\,n^{g}_{X,t}-\omega_X\,\mathcal{A}_{X,t}.
\label{eq:bank-div}
\end{equation}
Deposits are then the residual funding need, $\mathrm{dep}_{X,t}=\mathcal{A}_{X,t}-n_{X,t}$,
and the obligation state evolves by \eqref{eq:P-state}. Total dividends received by
households consolidate retail profits, capital-producer rents, and the banking sector:
\begin{equation}
\mathrm{Div}_{X,t}=(1-mc_X)\,Y_{X,t}+\Pi^{k}_{X,t}+\mathrm{div}^{\rm bank}_{X,t}
+\varsigma^{\rm reb}\,\zeta_X r^{wc}_{X,t}w_{X,t}N_{X,t},
\label{eq:dividends}
\end{equation}
where $\varsigma^{\rm reb}\in\{0,1\}$ is the working-capital rebate switch discussed
in Section~\ref{sec:wc} (baseline: $\varsigma^{\rm reb}=0$).

\subsection{Steady-state calibration of the agency friction}
\label{sec:bankcal}

At the steady state the constraint binds with equality and the system
\eqref{eq:phi-recursion}--\eqref{eq:mu-closed} collapses to a scalar fixed point.
Given a target leverage $\bar\theta$ and a target credit spread
$\bar\sigma\equiv \bar r^{k}-\bar r$, the pair $(\lambda_X,\omega_X)$ is solved from
\begin{align}
\bar\varphi&=\frac{\bar\Omega\,(1+\bar r)}{1-\bar\mu},
&
\bar\Omega&=\beta^{\rm int}\bigl[f+(1-f)\bar\varphi\bigr],
&
\bar\mu&=\frac{\bar\Omega\,\bar\sigma\,\bar\theta}{\bar\varphi},
\label{eq:ss-fixedpoint}\\[4pt]
\lambda_X&=\frac{\bar\varphi}{\bar\theta},
&
\omega_X&=\frac{\mathcal{D}}{\bar\theta}-(1-f)\,\bar\sigma,
&
\mathcal{D}&\equiv1-(1-f)(1+\bar r),
\label{eq:ss-lambda-omega}
\end{align}
where $\beta^{\rm int}$ is the bankers' discount factor, set to $1/(1+\bar r)$ so that
$\beta^{\rm int}R=1$ at the steady state. The scalar fixed point
\eqref{eq:ss-fixedpoint} can \emph{fold}: for high spread/leverage combinations the
map has no low root and the calibration is infeasible. The solver selects the least
root, which is also the one value iteration from below converges to.

Two internal consistency conditions define the steady state of the bank block: the
level of net worth implied by the binding constraint,
$\bar n^{IC}=\lambda\bar{\mathcal{A}}/\bar\varphi$, must equal the level implied by the
accumulation identity \eqref{eq:nw-lom},
\begin{equation}
\bar n^{ACC}
=\frac{1}{\mathcal{D}}
 \sum_{j\in\mathcal{J}}\Bigl[(1-f)\bigl(\bar R_{j}-\bar R\bigr)+\omega_X\Bigr]\bar a_j ,
\label{eq:ss-nw}
\end{equation}
and this equality is what pins the steady-state return on capital $\bar r^{k}_X$
in the first stage of the steady-state solve.

%======================================================================
\section{Sovereign risk and the default event}
\label{sec:sovrisk}
%======================================================================

\subsection{Debt instrument}

Both governments issue Hatchondo--Martinez/Chatterjee--Eyigungor perpetuities: a unit
of the stock pays a coupon $\delta_b$ this period and $(1-\delta_b)$ units of the same
bond survive into the next. Duration is approximately $1/\delta_b$ quarters. The
ex-coupon payoff to a unit of stock is $\Xi_{t+1}$ of equation~\eqref{eq:bond-payoff}.

\impl{Long duration is essential, not cosmetic. Because the price
$q_t$ capitalises the entire expected future path of default risk, a persistent
increase in $\pi^{d}$ produces a large mark-to-market loss on bank balance sheets.
Shortening duration to $\delta_b=0.25$ cuts the repricing by roughly a factor of six
and reverses the sign of the output response.}

\subsection{The exogenous, priced risk process}

Following Bocola (2016, eqs.~11--12), default risk is \textbf{exogenous}: it is a
priced probability, never a function of the debt stock or of any endogenous state.
A latent risk factor $s_t$ follows
\begin{equation}
s_{t+1}=(1-\rho_s)\,\bar s+\rho_s\,s_t+\sigma_s\,\varepsilon_{t+1},
\qquad \varepsilon_{t+1}\sim N(0,1),
\label{eq:s-process}
\end{equation}
and the one-quarter-ahead priced probability of default is the logistic transform
\begin{equation}
\pi^{d}_t\;=\;\pi^{d}(s_t)\;=\;\frac{1}{1+e^{-s_t}},
\qquad \bar s=\log\frac{0.001}{0.999}\ \ \text{so that}\ \ \pi^{d}(\bar s)=0.1\%.
\label{eq:pd}
\end{equation}

Two distinct objects must be kept separate:
\begin{itemize}
\item the \textbf{priced} probability $\pi^{d}_t$, which enters bond pricing
      \eqref{eq:bond-pricing} and every expected-return condition
      \eqref{eq:foc-general};
\item the \textbf{realised} indicator $d_t\in\{0,1\}$, which enters realised payoffs
      \eqref{eq:asset-payoff} and the government's flow budget.
\end{itemize}
The baseline experiment sets $\pi^{d}_t>0$ with $d_t\equiv0$: risk is priced but
never realised. This is precisely Bocola's ``pass-through'' design --- the entire
contraction is generated by the \emph{anticipation} of an event that does not occur.

\subsection{The default event}

The feared event is a single deterministic pure haircut: a write-down of the whole
$\D$-claim to a recovery value $\varrho_\D$, applied to coupon and continuation value
alike, so the survival factor is
\begin{equation}
h_t=1-d_t\,(1-\varrho_\D),
\qquad \varrho_\D=0.45 \ \ \text{(the 2012 Greek PSI; Bocola's $D=0.55$)} .
\label{eq:haircut}
\end{equation}
There are no exogenous scarring costs: no output cost of default, no capital-quality
shock, no bank recapitalisation. The recession in the default state arises
\emph{endogenously}, through the same balance-sheet mechanism
\eqref{eq:ng}--\eqref{eq:mu-closed} operating on a smaller asset base. The
$\F$-sovereign never defaults, so $\F$-bonds are safe in both states.

\subsection{The expectation operator}

The conditional expectations in \eqref{eq:omega}, \eqref{eq:foc-general},
\eqref{eq:bond-pricing} and \eqref{eq:agg-euler} are genuine multi-dimensional
integrals over (i) the innovation to $s$ and (ii) the discrete default fork. Let
$\{\varepsilon_m,w_m\}_{m=1}^{M}$ be the Gauss--Hermite nodes and weights of the
probabilists' Hermite rule, so that
$s_{t+1}^{(m)}=(1-\rho_s)\bar s+\rho_s s_t+\sigma_s\varepsilon_m$, and let
$d'\in\{0,1\}$ index the regime. For any function $g$ of next period's state,
\begin{equation}
\boxed{\;
\E_t\bigl[g\bigr]
=\sum_{m=1}^{M}w_m\Bigl\{
   \bigl(1-\pi^{d}_t\bigr)\,g\bigl(0,\mathcal{S}^{(m)}_{t+1}\bigr)
   \;+\;\pi^{d}_t\,\Bigl[
      \varkappa_{\rm tpi}\,g^{\rm hon}\bigl(\mathcal{S}^{(m)}_{t+1}\bigr)
      +\bigl(1-\varkappa_{\rm tpi}\bigr)\,g\bigl(1,\mathcal{S}^{(m)}_{t+1}\bigr)
   \Bigr]\Bigr\}\;}
\label{eq:expectation}
\end{equation}
where $g(d',\cdot)$ evaluates the equilibrium decision rules of regime $d'$ at a
\emph{reachable} next-period state --- never a frozen stand-in economy --- and
$g^{\rm hon}$ is the backstop-honoured branch of Section~\ref{sec:tpi}. Setting
$\varkappa_{\rm tpi}=0$ collapses \eqref{eq:expectation} to the two-branch
(no-backstop) case, and $\pi^{d}\equiv0$ nests the risk-neutral model exactly.

\subsection{Branch stochastic discount factors and the endogenous premium}

The banker discounts with the household kernel of its own country, evaluated
branch-by-branch on the GHH composite:
\begin{equation}
\Lambda^{(d')}_{X,t,t+1}
=\beta^{\rm int}_X
 \left(\frac{x_{X,t}}{x^{(d')}_{X,t+1}}\right)^{\sigma_X},
\qquad
\Omega^{(d')}_{X,t,t+1}
=\Lambda^{(d')}_{X,t,t+1}\Bigl[f+(1-f)\varphi^{(d')}_{X,t+1}\Bigr].
\label{eq:branch-sdf}
\end{equation}
Because consumption is lower in the default branch, $x^{(1)}<x^{(0)}$ and hence
$\Omega^{(1)}>\Omega^{(0)}$: the bank values wealth more in the bad state. This is
what makes the bond carry a genuine \emph{risk premium} rather than a pure
actuarial discount, and it is what makes $\E_t[\Omega]$ in \eqref{eq:mu-closed}
\emph{fall} when risk rises --- tightening the constraint.

\impl{A constant discount factor on the no-default branch (i.e.
$\Lambda^{(0)}\equiv\beta^{\rm int}$) cannot fall when risk rises, so
$\E_t[\Omega]$ climbs with $\varphi'$ and the constraint counterfactually
\emph{loosens}. The state-contingent kernel \eqref{eq:branch-sdf} is required for
the constraint to tighten with risk.}

The bond price \eqref{eq:bond-pricing} then decomposes into three pieces:
\begin{equation}
q^{\D}_t
=\underbrace{\frac{\E_t[\Xi^{\D}_{t+1}]}{R_{\D,t}}}_{\text{expected payoff}}
\;+\;\underbrace{\frac{\mathrm{Cov}_t\bigl(\Omega,\Xi^{\D}_{t+1}\bigr)}
                       {\E_t[\Omega]R_{\D,t}}}_{\text{risk premium}\;(<0)}
\;-\;\underbrace{\frac{\lambda_\D\mu_{\D,t}}
     {\E_t[\Omega]R_{\D,t}+\lambda_\D\mu_{\D,t}}\;q^{\D}_t}_{\text{liquidity/constraint discount}} .
\label{eq:bond-decomposition}
\end{equation}

%======================================================================
\section{Government and policy}
%======================================================================

\subsection{The consolidated government budget constraint}

The government of country $X$ finances an exogenous expenditure stream $G_X$, coupon
payments on the surviving stock, and any bailout outlay, with lump-sum taxes and new
issuance. In own-good units,
\begin{equation}
\boxed{\;
G_X+\delta_b\,B_{X,t}\,h_t
\;=\;T^{\tau}_{X,t}
\;+\;q^{X}_{t}\Bigl[\,B_{X,t+1}-(1-\delta_b)\,B_{X,t}\,h_t\Bigr]\;}
\label{eq:govt-budget}
\end{equation}
so that the debt stock evolves as
\begin{equation}
B_{X,t+1}=(1-\delta_b)\,B_{X,t}\,h_t
\;+\;\frac{G_X+\delta_b B_{X,t}h_t-T^{\tau}_{X,t}}{q^{X}_{t}} .
\label{eq:debt-lom}
\end{equation}
Because $q^{X}_t$ falls when risk rises, \eqref{eq:debt-lom} contains the standard
\emph{financing} amplification: the same primary deficit requires more face value to
be placed when the sovereign is under stress. The end-of-period stock is
forward-integrated \emph{inside} every equilibrium evaluation and absorbed by banks;
clearing against a fixed $\bar B$ instead opens a Walras leak of order $0.5\%$ of GDP
per $5\%$ debt deviation.

\subsection{The fiscal rule}

Taxes follow a Bohn (1998) rule in the form Bocola estimates, i.e. as a constant
\emph{elasticity} of the tax level with respect to the surviving debt stock:
\begin{equation}
\boxed{\;
T^{\tau}_{X,t}=\bar T^{\tau}_X
\left(\frac{B_{X,t}\,h_t}{\bar B_X}\right)^{\gamma_\tau},
\qquad \gamma_\tau=1 \;}
\label{eq:bohn}
\end{equation}
with the steady-state anchor $\bar T^{\tau}_X=G_X+\delta_b\bar B_X(1-\bar q^{X})$
obtained from \eqref{eq:govt-budget} at a constant stock.

\begin{remark}[Elasticity, not level]
Reading $\gamma_\tau$ as a coefficient on the debt \emph{level} rather than as an
elasticity turns a $55\%$ haircut into a fiscal windfall on the order of $50\%$ of
GDP, which makes the default event expansionary and renders most of the $d=1$ region
unsolvable. Under \eqref{eq:bohn} a haircut to $\varrho_\D=0.45$ delivers
$B'/\bar B=0.451$, as it should: the fiscal-relief leg of default is real but
second-order relative to the bank-loss leg at the calibrated exposure.
\end{remark}

\subsection{The monetary block of the union}
\label{sec:money}

Prices are fully flexible and there is no nominal rigidity, so the model contains no
Taylor rule and no nominal interest rate to set. What a monetary union \emph{does}
impose on real allocations is modelled directly, in two parts.

\paragraph{(i) A single union-wide funding market.} The two national deposit-market
clearing conditions are replaced by one union-wide clearing condition (in $\D$-good
units) together with a no-arbitrage condition across national deposit legs. Deposits
are own-good claims remunerated at national rates, and a frictionless union interbank
market equalises their real returns:
\begin{equation}
\boxed{\;
\bigl(1+r_{\D,t}\bigr)
\;=\;\bigl(1+r_{\F,t}\bigr)\,\frac{\E_t\bigl[p_{t+1}\bigr]}{p_t}
\;-\;\kappa_{\rm nfa}\,\frac{P_{\D,t}-P_{\F,t}}{\bar Y_\D}\;}
\label{eq:uip}
\end{equation}
This \textbf{deposit-UIP} condition is the flexible-price image of a single nominal
policy rate plus national inflation differentials, i.e. the Backus--Kehoe--Kydland /
Baxter--Crucini single-traded-bond margin. Together with \eqref{eq:agg-euler} for
$\D$ it determines both national deposit rates. Because pass-through through the
union interbank is zero-profit, the cross-border deposit position is the absorption
margin that breaks the national saving-equals-investment trap; imposing the literal
condition $r_{\D,t}=r_{\F,t}$ on own-good legs instead is \emph{incorrect}, because
it leaves an unassigned real-exchange-rate valuation profit and opens a Walras leak.

\paragraph{(ii) A stationarity-inducing external premium.} The last term of
\eqref{eq:uip} is the Schmitt-Grohé--Uribe debt-elastic premium on the net external
position (Bocola's only foreign friction), proxied here by the cross-country wealth
imbalance $P_{\D,t}-P_{\F,t}$. It is exactly zero at the symmetric steady state, so
it is undistorting there and only induces stationarity off-steady-state.
$\kappa_{\rm nfa}=0$ nests frictionless parity.

\subsection{The TPI/OMT backstop}
\label{sec:tpi}

% STALE AS OF 2026-08-30. This section describes the earlier three-branch design, in
% which the DEFAULT fork split into backstop-honoured and reneged and the central bank
% bought nothing. The implemented policy is now a ONE-SIDED YIELD PEG WITH REAL
% PURCHASES: with per-period probability phi the central bank is in the D-sovereign
% market and buys until the price reaches Q*, it holds what it buys (a state, running
% off at delta_b), it is funded by household claims rather than bank reserves, and it
% remits the carry to the treasury. Rewrite this section from point_map.py before the
% draft goes out.


The union's central bank operates a sovereign backstop of the OMT/TPI type. Two
formulations are implemented; both are \emph{announcement}, not rate, instruments.

\paragraph{(a) Priced activation (the production formulation).} The default fork of
\eqref{eq:expectation} splits into a \textbf{backstop-honoured} and a
\textbf{backstop-reneged} branch, weighted by the priced activation probability
$\varkappa_{\rm tpi}\in[0,1]$. In the honoured branch the central bank (i) redeems
the $\D$-bond at a near-par value $\varrho^{\rm tpi}$ and (ii) averts a fraction
$\varsigma$ of the recession, so that the honoured continuation is a blend of the two
regimes' decision rules:
\begin{align}
g^{\rm hon}(\mathcal{S}')&=(1-\varsigma)\,g\bigl(1,\mathcal{S}'\bigr)
                          +\varsigma\,g\bigl(0,\mathcal{S}'\bigr),
\label{eq:tpi-blend}\\[3pt]
\Xi^{\D,\rm hon}_{t+1}&=\varrho^{\rm tpi}\Bigl[\delta_b+(1-\delta_b)q^{\D,\rm hon}_{t+1}\Bigr],
\qquad
\Xi^{\D,\rm ren}_{t+1}=\varrho_\D\Bigl[\delta_b+(1-\delta_b)q^{\D,(1)}_{t+1}\Bigr].
\label{eq:tpi-payoffs}
\end{align}
At $\varrho^{\rm tpi}=\varsigma=1$ a honoured backstop fully neutralises the event,
so the effective default probability becomes $\pi^{d}_t(1-\varkappa_{\rm tpi})$ ---
the textbook reading that a fully credible backstop removes the risk premium at zero
purchases. At $\varkappa_{\rm tpi}=0$ or $\varsigma=0$ the two-branch solve is nested
exactly.

\paragraph{(b) A mechanical price floor (the Markov-switching formulation).} In the
regime $\mathfrak{s}^{\rm tpi}_t=1$ the central bank supports the $\D$-bond at a
floor that prices the \emph{same} expected payoff at the steady-state constraint
spread, purchasing quantity $\mathcal{Q}^{cb}_t$ according to a portfolio-balance rule:
\begin{equation}
q^{\rm floor}_t=\frac{\E_t\bigl[\Xi^{\D}_{t+1}\bigr]}{1+r_{\D,t}+\bar\lambda\bar\mu/\bar\Omega},
\qquad
\mathcal{Q}^{cb}_t=\frac{\max\bigl\{0,\;q^{\rm floor}_t-q^{\D,\rm free}_t\bigr\}}{\psi_{cb}}
\,\mathfrak{s}^{\rm tpi}_t,
\qquad
q^{\D}_t=q^{\D,\rm free}_t+\psi_{cb}\,\mathcal{Q}^{cb}_t .
\label{eq:tpi-floor}
\end{equation}
The floor strips out the liquidity component of the yield while leaving default
compensation intact. Any central-bank profit or loss is rebated to the fiscal
authority, so \eqref{eq:govt-budget} continues to hold as a consolidated constraint.
The $\max$ operator in \eqref{eq:tpi-floor} is smoothed in implementation, since it
is non-differentiable exactly where the floor switches on.

\subsection{Macroprudential policy}

The model contains \textbf{no macroprudential rule}: this is a deliberate scope
decision, not an omission of an implemented block. The natural policy object, were
one to be added, is the divertability parameter $\lambda_X$ of \eqref{eq:IC},
equivalently the leverage ceiling
\begin{equation}
\theta_{X,t}\;\le\;\bar\theta_{X,t}\;=\;\frac{\varphi_{X,t}}{\lambda_X},
\label{eq:macropru}
\end{equation}
which a capital requirement would tighten. The comparative statics of such a policy
are fully characterised in closed form by \eqref{eq:ss-fixedpoint}: a permanently
tighter $\lambda$ lowers steady-state leverage one-for-one and, through
$\bar\mu=\bar\Omega\bar\sigma\bar\theta/\bar\varphi$, raises the steady-state credit
spread \eqref{eq:rwc} --- trading a smaller crisis amplification for a permanently
higher cost of capital. Making $\lambda$ state-contingent (e.g.\ a countercyclical
buffer $\lambda_t=\bar\lambda\,\exp\{-\phi_{\rm mp}\log(n_{X,t}/\bar n_X)\}$) would
enter the system only through \eqref{eq:mu-closed} and would require no other change.

%======================================================================
\section{Equilibrium}
\label{sec:equilibrium}
%======================================================================

\subsection{Market clearing}

\paragraph{Goods markets.} Each country's good clears against domestic absorption
and net exports,
\begin{align}
Y_{\D,t}&=P^{c}_{\D,t}C_{\D,t}+I_{\D,t}+G_\D+NX_{\D,t},
\label{eq:goods-D}\\
Y_{\F,t}&=P^{c}_{\F,t}C_{\F,t}+I_{\F,t}+G_\F+NX_{\F,t}.
\label{eq:goods-F}
\end{align}
The term $P^{c}_{X,t}C_{X,t}$ is consumption expenditure valued in own-good units;
subtracting imports and adding exports through $NX$ recovers the own-good content.
By Walras' law, \eqref{eq:goods-F} and the current-account identity are implied by
the remaining conditions and are therefore \emph{dropped} from the solved system and
monitored as diagnostics.

\paragraph{Deposit market (union-wide).} Household saving equals bank funding, by
quantity, with the union interbank reallocating across countries:
\begin{equation}
A_{\D,t}P^{c}_{\D,t}+p_t\,A_{\F,t}P^{c}_{\F,t}
=\mathrm{dep}_{\D,t}+p_t\,\mathrm{dep}_{\F,t},
\qquad
A_{X,t}=\frac{\mathrm{dep}_{X,t}}{P^{c}_{X,t}} ,
\label{eq:deposit-clearing}
\end{equation}
with the cross-border deposit position $\mathrm{nfa}^{\rm dep}_{\D,t}$ as the
absorption margin, and the rate structure pinned by \eqref{eq:agg-euler} and
\eqref{eq:uip}.

\paragraph{Sovereign bond markets.} Banks absorb the entire outstanding stock of each
sovereign:
\begin{equation}
b^{\D}_{\D,t+1}+b^{\D}_{\F,t+1}=B_{\D,t+1},
\qquad
b^{\F}_{\F,t+1}+b^{\F}_{\D,t+1}=B_{\F,t+1}.
\label{eq:bond-clearing}
\end{equation}
The cross-border split is held at its calibrated share in the projection solver
($b^{\D}_{\F}/B_\D$ fixed at the EBA-based exposure); in the path-based
implementation it is instead determined by the cross-border FOC with quadratic
portfolio adjustment costs,
\begin{equation}
b^{\F}_{\D,t}=\bar b^{\F}_{\D}
+\frac{1}{\psi^{b}}\Bigl[\E_t\bigl[R^{\F}_{t+1}\bigr]-R_{\D,t}
-\overline{\mathrm{exc}}-\lambda_\D\mu_{\D,t}/\E_t[\Omega_\D]\Bigr].
\label{eq:xborder-foc}
\end{equation}

\paragraph{Capital and labour.} Capital market clearing is implicit in the
intermediary's balance sheet --- banks hold the entire capital stock, and the
capital Euler \eqref{eq:capital-euler} determines $K_{X,t+1}$. Labour clears at the
wage satisfying \eqref{eq:agg-labour} and \eqref{eq:labour-demand} jointly.

\paragraph{External account.} The current account identity
\begin{equation}
NX_{\D,t}+\bigl(\text{net cross-border factor income}\bigr)
=\Delta\bigl(\text{net foreign asset position}\bigr)
\label{eq:ca}
\end{equation}
holds as an implication and is monitored as a Walras diagnostic.

\subsection{Recursive competitive equilibrium}

\begin{definition}[Recursive competitive equilibrium]
\label{def:rce}
Let the aggregate state be
\begin{equation}
\mathcal{S}_t=\bigl(K_{\D,t},\,K_{\F,t},\,P_{\D,t},\,P_{\F,t},\,B_{\D,t},\,s_t,\,Z_{\D,t}\bigr)
\;\in\;\R^{7},
\qquad d_t\in\{0,1\},
\label{eq:state}
\end{equation}
i.e.\ the two capital stocks, the two banks' gross deposit obligations, the risky
debt stock, the sovereign-risk factor and TFP, together with the discrete default
regime. A recursive competitive equilibrium consists of
\begin{enumerate}
\item decision rules
$\bigl\{N_X(d,\mathcal{S}),\,K'_X(d,\mathcal{S}),\,r_X(d,\mathcal{S}),\,p(d,\mathcal{S})\bigr\}$;
\item valuation rules
$\bigl\{\varphi_X(d,\mathcal{S}),\,q^{X}(d,\mathcal{S})\bigr\}$
and household aggregates $\bigl\{C_X(d,\mathcal{S}),\,A_X(d,\mathcal{S})\bigr\}$;
\item multipliers $\mu_X(d,\mathcal{S})\ge0$ and a transition law
$\mathcal{S}'=\mathcal{T}\bigl(d,\mathcal{S},d',\varepsilon'\bigr)$;
\end{enumerate}
such that, at every $(d,\mathcal{S})$: households optimise
(\eqref{eq:agg-labour}, \eqref{eq:agg-euler}) subject to
\eqref{eq:hh-budget-agg}; firms optimise
(\eqref{eq:labour-demand}--\eqref{eq:mpk}) and capital producers satisfy
\eqref{eq:tobins-q}; intermediaries optimise
(\eqref{eq:foc-general}, \eqref{eq:phi-recursion}) subject to the incentive
constraint \eqref{eq:IC} with complementarity \eqref{eq:kkt}; the government
satisfies \eqref{eq:govt-budget}--\eqref{eq:bohn}; all markets clear
(\eqref{eq:goods-D}, \eqref{eq:deposit-clearing}, \eqref{eq:bond-clearing},
\eqref{eq:uip}); and expectations are formed under \eqref{eq:expectation} with
$\mathcal{S}'$ generated by the same rules.
\end{definition}

\subsection{The solved system, equation by equation}
\label{sec:system}

At each point $(d,\mathcal{S})$ the equilibrium is characterised by seven equations
in seven unknowns,
\begin{equation}
\mathbf{x}=\bigl(N_\D,\;N_\F,\;K'_\D,\;K'_\F,\;r_\D,\;r_\F,\;p\bigr),
\label{eq:unknowns}
\end{equation}
listed in Table~\ref{tab:system}. The valuations $\varphi_X$, $q^{X}$ and the
multiplier $\mu_X$ are \emph{read off} the recursions
\eqref{eq:phi-recursion}, \eqref{eq:bond-pricing} and \eqref{eq:mu-closed} given the
continuation, rather than solved for --- which is what breaks the simultaneity
between the bond price and next period's state.

\begin{table}[htbp]
\centering
\small
\caption{The seven market-clearing conditions and the unknown each pins down.}
\label{tab:system}
\begin{tabular}{@{}clll@{}}
\toprule
\# & Condition & Equation & Pins \\
\midrule
1 & Capital Euler, $\D$              & $\E[\Omega_\D(r^{k}_\D{}'-r_\D)]-\lambda_\D\mu_\D=0$ & $K'_\D$ \\
2 & Capital Euler, $\F$              & $\E[\Omega_\F(r^{k}_\F{}'-r_\F)]-\lambda_\F\mu_\F=0$ & $K'_\F$ \\
3 & Labour market, $\D$              & $\chi_\D N_\D^{1/\nu_\D}-w_\D/P^{c}_\D=0$            & $N_\D$   \\
4 & Labour market, $\F$              & $\chi_\F N_\F^{1/\nu_\F}-w_\F/P^{c}_\F=0$            & $N_\F$   \\
5 & Household deposit Euler, $\D$    & $x_\D^{-\sigma}-\tilde\beta_\D\E[(1+r^{c}_\D{}')x_\D'^{-\sigma}]=0$ & $r_\D$ \\
6 & Deposit-UIP (union)              & $(1+r_\D)-(1+r_\F)\E[p']/p+\kappa_{\rm nfa}\Delta P/\bar Y=0$ & $r_\F$ \\
7 & Goods market, $\D$               & $Y_\D-P^{c}_\D C_\D-I_\D-NX_\D-G_\D=0$               & $p$      \\
\midrule
\multicolumn{4}{@{}l@{}}{\emph{Dropped (Walras-redundant), monitored as diagnostics:} goods market $\F$; current account.}\\
\bottomrule
\end{tabular}
\end{table}

\subsection{Definitions of the key endogenous variables}

The objects reported in the experiments are defined as follows.

\begin{align}
\text{Credit / lending spread (annualised bp)}\quad
&\mathrm{sp}_{X,t}=4\times10^{4}\cdot\frac{\lambda_X\,\mu_{X,t}}{\varphi_{X,t}},
\label{eq:def-spread}\\[3pt]
\text{Bank leverage}\quad
&\theta_{X,t}=\frac{\mathcal{A}_{X,t}}{n_{X,t}}
   \;\le\;\frac{\varphi_{X,t}}{\lambda_X},
\label{eq:def-leverage}\\[3pt]
\text{Sovereign yield / spread over the safe rate}\quad
&y^{\D}_t=\frac{\delta_b+(1-\delta_b)q^{\D}_t}{q^{\D}_t}-1,
\qquad
\mathrm{spr}^{\D}_t=y^{\D}_t-y^{\F}_t,
\label{eq:def-yield}\\[3pt]
\text{Constraint slack}\quad
&\mathrm{slack}_{X,t}=\varphi_{X,t}n_{X,t}-\lambda_X\mathcal{A}_{X,t}\;\ge\;0,
\label{eq:def-slack}\\[3pt]
\text{Real exchange rate}\quad
&\mathrm{rer}_t=\frac{p_t\,P^{c}_{\F,t}}{P^{c}_{\D,t}},
\qquad
\text{Debt/GDP}\;=\;\frac{q^{\D}_tB_{\D,t}}{4Y_{\D,t}} .
\label{eq:def-rer}
\end{align}

\subsection{Decomposition of the output response}

Because production \eqref{eq:production} is Cobb--Douglas and capital is
predetermined, the impact response of output is an exact identity in three terms:
\begin{equation}
d\log Y_{\D,t}=d\log Z_{\D,t}+\alpha_\D\,d\log K_{\D,t}+(1-\alpha_\D)\,d\log N_{\D,t}.
\label{eq:decomp-Y}
\end{equation}
Substituting the labour-market condition \eqref{eq:agg-labour} together with
\eqref{eq:labour-demand} and solving for hours,
\begin{equation}
\boxed{\;
d\log N_{\D,t}
=\Theta\Bigl[\,d\log Z_{\D,t}+\alpha_\D\,d\log K_{\D,t}
-\!\!\underbrace{d\log\bigl(1+\zeta_\D r^{wc}_{\D,t}\bigr)}_{\text{financial wedge}}
\!\!-\,d\log P^{c}_{\D,t}\Bigr],
\quad
\Theta=\frac{1}{1/\nu_\D+\alpha_\D}\;}
\label{eq:decomp-N}
\end{equation}
and, using \eqref{eq:rwc}, the financial wedge splits additively into a
\textbf{deposit-rate} leg and a \textbf{credit-spread} leg,
$r^{wc}=r_{\D}+\lambda_\D\mu_\D/\E[\Omega_\D]$. The second is the Bocola channel:
it is the only term through which an increase in the priced probability of sovereign
default reaches output.

%======================================================================
\section{The steady state}
%======================================================================

The deterministic steady state is imposed to be \textbf{symmetric}: country
asymmetries enter through shocks only, so $p_{ss}=1$ and every real quantity is
common across countries. It is computed in two stages.

\paragraph{\texorpdfstring{Stage 1: $\{\bar r^{k}_\D,\bar r^{k}_\F,\bar p\}$.}
                          {Stage 1.}} Two capital-market
conditions --- the equality of $\bar n^{IC}$ and $\bar n^{ACC}$ from
Section~\ref{sec:bankcal}, one per country --- plus external balance,
\begin{equation}
NX_\D+\bigl(\text{net cross-border coupon income}\bigr)=0 ,
\label{eq:ss-external}
\end{equation}
determine the two returns on capital and the terms of trade. TFP is then rescaled to
normalise $\bar Y_X=1$ and $\bar N_X=1$, and $\chi_X$ is set from the GHH labour
condition \eqref{eq:agg-labour}.

\paragraph{\texorpdfstring{Stage 2: $\{\beta_\D,\beta_\F\}$.}
                          {Stage 2.}} The heterogeneous-agent block is solved
for each candidate discount factor and the deposit-market condition
\begin{equation}
A^{ss}_X\bigl(\beta_X\bigr)=\bigl(\bar\theta_X-1\bigr)\bar n_X
\label{eq:ss-deposit}
\end{equation}
is inverted for $\beta_X$. Under the symmetric-steady-state doctrine
$\beta_\F=\beta_\D$ and the $\F$ condition is verified rather than imposed --- at a
symmetric steady state the union clearing coincides with each national market and the
cross-border deposit position is zero.

At the calibrated steady state the incentive constraint \textbf{binds}
($\bar\mu\approx0.02$, $\mathrm{slack}=0$), which is a deliberate choice: Bocola's own
posterior puts $\bar\mu\approx0.001$, i.e.\ exactly on the KKT kink, where any shock
drives some period's multiplier across zero and the solver loses a valid Jacobian.
The Gertler--Karadi (2011) robustly-binding calibration (leverage 4, a 200\,bp annual
spread) keeps the constraint strictly binding through both experiments at the cost of
smaller --- but correctly signed --- amplification.

%======================================================================
\section{Numerical solution}
%======================================================================

The model is solved \textbf{globally}, as recursive decision rules over the seven
continuous states \eqref{eq:state} and the discrete regime $d$, with separate
coefficient sets per regime.

\paragraph{Approximation.} Every rule is a Chebyshev polynomial on a Smolyak sparse
grid built from nested Chebyshev-extrema point sets (Krueger--Kubler): the grid is
the union of tensor products of per-level ``new point'' sets over multi-indices $i$
with $\sum_j(i_j-1)\le\ell$, and the basis uses the same index set over ``new
degree'' sets, so the collocation matrix is square and interpolation is exact at the
nodes. Anisotropic levels $\ell_j$ concentrate resolution on the risk factor $s$ and
the two net-worth states, where the constraint boundary moves. Rules that must remain
positive are collocated in logs; interest rates are carried as gross rates $1+r$.

\paragraph{Algorithm (time iteration).} Each sweep freezes the previous iterate as
the continuation, solves the seven residuals of Table~\ref{tab:system} pointwise by
Newton's method at every grid node and in both regimes, reads the valuations off
\eqref{eq:phi-recursion}--\eqref{eq:mu-closed}, then damps and refits the
coefficients. Convergence requires \emph{both} a settled rule and every point
clearing; nodes that fail to clear retain their previous values and are
down-weighted in the fit, so a failed corner cannot poison the global basis.

\paragraph{Smoothing of kinks.} Guards that would otherwise plant a plateau inside a
fitted object --- the caps on $\mu$ and $\varphi$, the net-worth floor of
\eqref{eq:nw-lom}, the $\max$ in \eqref{eq:tpi-floor} --- are replaced by the smooth
counterparts $\mathrm{smax}(x,\underline{x})=\underline{x}+\tfrac12[(x-\underline{x})
+\sqrt{(x-\underline{x})^2+\epsilon^2}]$. The KKT switch in \eqref{eq:mu-closed} is
\emph{not} smoothed: it is the complementarity itself, not a numerical guard.

\paragraph{Known limitations.} The magnitudes reported at the baseline are indicative
rather than estimated: the solve converges at approximation level $\ell=1$; $\ell=2$
does not converge, because the near-unit-root corners of the $d'=1$ regime poison the
global basis. The heterogeneous-agent distribution is carried by its mean
(Remark~\ref{rem:closure}), the household deposit Euler does not price the default
branch (households' deposits are riskless, faithful to Bocola), and the cross-border
bond split is fixed at its calibrated share in the projection solver.

%======================================================================
\section{Calibration}
%======================================================================

\begin{table}[htbp]
\centering
\small
\caption{Baseline quarterly calibration. Symmetric across countries unless noted.}
\label{tab:calibration}
\begin{tabular}{@{}llll@{}}
\toprule
Parameter & Symbol & Value & Source / target \\
\midrule
\multicolumn{4}{@{}l}{\emph{Households}}\\
Inverse EIS                       & $\sigma$        & $1.0$    & Bocola \S II.A.1 (log utility)\\
Frisch elasticity                 & $\nu$           & $2.0$    & Bocola Table 1 ($1/\nu=0.5$)\\
Persistence of idiosyncratic risk & $\rho_e$        & $0.90$   & standard\\
S.d.\ of idiosyncratic risk       & $\sigma_e$      & $0.20$   & standard\\
Borrowing limit                   & $\underline{a}$ & $0$      & no borrowing\\
\midrule
\multicolumn{4}{@{}l}{\emph{Production and capital}}\\
Capital share                     & $\alpha$        & $0.30$   & Bocola Table 1\\
Depreciation                      & $\delta$        & $0.025$  & standard\\
Elasticity of $q$ to $I/K$        & $\xi$           & $0.42$   & Bocola Table 2 (posterior mean)\\
Demand elasticity                 & $\epsilon$      & $6.0$    & $20\%$ markup\\
Working-capital share             & $\zeta$         & $1.0$    & Neumeyer--Perri; Bocola \S V.C\\
\midrule
\multicolumn{4}{@{}l}{\emph{Intermediaries}}\\
Exit/payout share                 & $f$             & $0.04$   & Bocola Table 2 ($\psi=0.96$ survival)\\
Bankers' discount factor          & $\beta^{\rm int}$ & $0.997$& $=1/R^{bg}$, Bocola Table 1\\
Steady-state deposit rate         & $\bar r$        & $0.003$  & sample-average risk-free rate\\
Target leverage                   & $\bar\theta$    & $4.0$    & Gertler--Karadi (2011)\\
Target credit spread              & $\bar\sigma$    & $0.005$  & $200$\,bp p.a.\ (GK11)\\
Implied multiplier                & $\bar\mu$       & $\approx0.02$ & strictly binding\\
\midrule
\multicolumn{4}{@{}l}{\emph{Sovereign debt and risk}}\\
Amortisation rate                 & $\delta_b$      & $0.056$  & Bocola Table 1; duration $\approx$ 7\,y\\
Debt stock                        & $\bar B$        & $0.98$   & bank exposure $=7.6\%$ of assets\\
Recovery rate                     & $\varrho_\D$    & $0.45$   & Greek PSI 2012\\
Persistence of $s$                & $\rho_s$        & $0.95$   & Bocola Table 2\\
Innovation s.d.\ of $s$           & $\sigma_s$      & $0.63$   & Bocola Table 2\\
Rest-point default probability    & $\pi^{d}(\bar s)$ & $0.1\%$ & normalisation\\
\midrule
\multicolumn{4}{@{}l}{\emph{Government, trade and policy}}\\
Fiscal elasticity                 & $\gamma_\tau$   & $1.0$    & Bocola Table 1\\
Government spending               & $G$             & $0$      & normalisation\\
Home bias                         & $\varpi$        & $0.85$   & standard\\
Trade elasticity                  & $\eta$          & $0.5$    & standard\\
External premium                  & $\kappa_{\rm nfa}$ & $0$   & frictionless UIP nested\\
TPI redemption value              & $\varrho^{\rm tpi}$ & $1.0$ & near-par floor\\
TPI real shield                   & $\varsigma$     & $0.5$    & partial backstop\\
TPI activation                    & $\varkappa_{\rm tpi}$ & $\{0,\,0.5,\,1\}$ & per experiment\\
CB portfolio-balance elasticity   & $\psi_{cb}$     & $0.5$    & solver-feasible floor\\
\bottomrule
\end{tabular}
\end{table}

%======================================================================
\appendix
\section{Symbol-to-code map}
%======================================================================

\begin{table}[htbp]
\centering
\small
\caption{Mapping between the notation of this document and \code{code/global/}.}
\label{tab:codemap}
\begin{tabular}{@{}l>{\raggedright\arraybackslash}p{5.3cm}l@{}}
\toprule
Symbol & Code name & Location \\
\midrule
$\varphi_{X,t}$ (franchise value)      & \code{alpha\_D}, \code{alpha\_F}        & \code{point\_map.py}, \code{bank.py}\\
$\mu_{X,t}$ (IC multiplier)            & \code{mu\_D}, \code{mu\_F}              & \code{point\_map.py}\\
$\lambda_X$ (divertability)            & \code{lambda\_K}, \code{lambda\_bD}, \code{lambda\_bF} & \code{calibration.py}\\
$\Omega_{X,t,t+1}$                     & \code{Om0\_D}, \code{Om1\_D}, \code{E\_Om\_D} & \code{point\_map.py}\\
$\mathcal{A}_{X,t}$ (bank assets)      & \code{assets\_D}, \code{assets\_F}      & \code{point\_map.py}\\
$\lambda\mathcal{A}$ (divertable base) & \code{lev\_D}, \code{lev\_F}            & \code{point\_map.py}\\
$n^{g}_{X,t}$, $n_{X,t}$               & \code{ng\_D}, \code{n\_D}               & \code{point\_map.py}\\
$P_{X,t}$ (deposit obligation state)   & \code{P\_D}, \code{Pp\_D}               & \code{point\_map.py}, \code{state\_grid.py}\\
$P^{c}_{X,t}$ (CES price index)        & \code{P\_CES\_D}                        & \code{trade.py}\\
$q^{\D}_t$, $q^{\F}_t$                 & \code{Q\_bD}, \code{Q\_bF}              & \code{point\_map.py}\\
$Q_{X,t}$ (Tobin's $q$)                & \code{Q\_D}, \code{Q\_F}                & \code{capital.py}\\
$h_t$ (survival factor)                & \code{surv}, \code{surv\_d}             & \code{point\_map.py}\\
$\pi^{d}(s)$                           & \code{default\_prob(s)}                 & \code{state\_grid.py}\\
$r^{wc}_{X,t}$                         & \code{r\_wc\_D}, \code{wedge\_sp\_D}    & \code{point\_map.py}\\
$\varkappa_{\rm tpi},\varrho^{\rm tpi},\varsigma$ & \code{tpi\_activation}, \code{recovery\_tpi\_D}, \code{tpi\_real\_shield} & \code{calibration.py}\\[2pt]
$\bar T_X$ (accounting anchor)         & \code{hh\_T\_D}, \code{hh\_T\_F}        & \code{recursive\_main.py}\\
$\tilde\beta_X$                        & \code{beff\_D}, \code{beff\_F}          & \code{point\_map.py}\\
$\mathcal{S}_t$ (state vector)         & \code{STATE\_NAMES}                     & \code{state\_grid.py}\\
$\mathbf{x}$ (seven unknowns)          & \code{SOLVE7}                           & \code{decision\_rules.py}\\
\bottomrule
\end{tabular}
\end{table}

\section{Summary of the transmission mechanism}

For reference, the chain of equations through which an exogenous increase in the
priced probability of sovereign default contracts real activity:
\begin{enumerate}
\item $s_t\uparrow$ raises $\pi^{d}_t$ through \eqref{eq:pd};
\item the multi-branch expectation \eqref{eq:expectation} places more weight on the
      haircut payoff $\Xi^{\D,\rm ren}$, so $q^{\D}_t$ falls through
      \eqref{eq:bond-pricing} --- amplified by duration $1/\delta_b$;
\item bank gross wealth $n^{g}_{\D}=\mathcal{X}_\D-P_\D$ falls in \eqref{eq:ng}
      (mark-to-market loss), so $n_\D$ falls in \eqref{eq:nw-lom};
\item simultaneously $\E_t[\Omega_\D]$ falls, because the state-contingent kernel
      \eqref{eq:branch-sdf} puts weight on a low-consumption branch;
\item both effects raise $\mu_{\D,t}$ in the closed form \eqref{eq:mu-closed}: the
      incentive constraint tightens;
\item the credit spread $\lambda_\D\mu_\D/\E[\Omega_\D]$ widens \eqref{eq:rwc},
      raising the cost of working capital;
\item the effective wage rises in \eqref{eq:labour-demand}, hours fall through
      \eqref{eq:agg-labour}, and output falls by \eqref{eq:decomp-N};
\item $\mu_\D\uparrow$ also raises the required return on capital
      \eqref{eq:capital-euler}, so $K'_\D$ and investment fall, propagating the
      contraction into subsequent periods;
\item the government must place more face value at a lower price
      \eqref{eq:debt-lom}, taxes rise through \eqref{eq:bohn}, and the loop repeats.
\end{enumerate}
No default is ever realised: $d_t\equiv0$ throughout. The entire contraction is the
pass-through of \emph{priced} risk.

\end{document}
